\subsubsection{Scheme A (minimal): composite vertices \texorpdfstring{$W_n$}{Wn} on the
smallest tower}
\label{ped:sec:allorders:min}

Use the fewest fields ($\Sig_k$ first at $n=3(k{-}1)+1$, i.e.\ $k\le\lceil n/3\rceil$,
by the all-orders counting) and put every subtraction \emph{inside the vertex kernel}: at order $n$
add the single coupling
\begin{equation}
	\frac{\gs^{n-2}}{n!}\,W_n(\Sig_a,\Sig_b,\Sig_c),
	\qquad n=a+b+c,\quad\text{balanced partition } a,b,c\le\lceil n/3\rceil,
	\label{ped:eq:schemeA}
\end{equation}
with the recursion established through $n=6$:
\begin{equation}
	W_4=V_4,\qquad W_5=V_5,\qquad
	W_6=V_6-\tfrac65\,\big(W_5\circ V_4^{-1}\circ W_5\big),\qquad\dots
	\label{ped:eq:Wnrecursion}
\end{equation}
--- in general $W_n=V_n-(\text{all colour-2-reducible composites of lower }W_m)$, the
coloured-operad statement that the coloured vertex equals the amplitude minus the exchange
regions, with the $\Sig$-propagators moduli-free kernels. A reading
of why this scheme is natural: the subtraction removes from the
elementary vertex region exactly the part that can be \emph{covered} by lower-order
auxiliary-field exchanges --- mirroring how the original closed-SFT action is itself
constructed, by identifying the gaps in the moduli-space covering left by propagator
regions and covering them with vertex regions; here the \enquote{propagators} are the
moduli-free colour-2 necks and the covering bookkeeping happens at the level of kernels
rather than regions. Proven:
$n\le6$ (\cref{p:sec:hs:quintic,p:sec:hs:sextic}). At $n=7$ the literal one-coupling rule
\eqref{ped:eq:schemeA} is now \emph{falsified} (\cref{p:sec:hs:septic}):
$\Sig_3$ becomes unavoidable ($7=3+2+2$ is the only partition reaching seven strands),
but a $(3,2,2)$ coupling alone cannot supply the $V_7$ contact --- with no $O(\gs^4)$
source the $\Sig_3$ saddle is $O(\gs)$ and the coupling's on-shell footprint is
$O(\gs^6)$, for \emph{any} kernel and coefficient. The $\Sig_3$
Hubbard--Stratonovich pair at $O(\gs^4)$ (scheme B's even pair) is \emph{forced}, so at
$n=7$ scheme A must restructure its $n=6$ sector and merge into the scheme-B family
(\cref{p:sec:hs:septic}).

\subsubsection{Scheme B (dressed, uniform): pair each field with the running effective
vertex}
\label{ped:sec:allorders:dressed}

The preferred bookkeeping (its first instance, at the sextic order, is proven,
\cref{p:sec:hs:sextic}). Define the \emph{running effective vertex}
\begin{equation}
	V_{n,\rm eff}
	:= V_n+\big(\text{all feedback composites generated by the orders}<n\big)
	\label{ped:eq:runningV}
\end{equation}
--- operationally, $V_{n,\rm eff}$ is the $n$-th derivative, on the Veronese variety, of
the action obtained by eliminating the auxiliaries of the orders $<n$; at $n=6$ this is
$\tilde V_6$ with $\tilde V_6(P_3,P_3)=W_6(\Psi^6)$, built by the $\Phi$-map
\eqref{ped:eq:tildeMdef}. The scheme then adds, for each $k\ge2$, one uniform triple of
terms --- an HS pair (kinetic $+$ source, ratio $1{:}2$) at the even order $n=2k$ and
one odd-order coupling at $n=2k+1$:
\begin{equation}
	\boxed{\;
	S^{\rm dr}
	= \tfrac12V_2(\Psi^2)+\frac{\gs}{3!}V_3(\Psi^3)
	+\sum_{k\ge2}\bigg[
	\frac{\gs^{2k-2}}{(2k)!}\Big(-\tilde M_k(\Sig_k,\Sig_k)
	+2\,\Mean{\tilde u_k(\Sig_{k-1}),\Sig_k}\Big)
	+\frac{\gs^{2k-1}}{(2k+1)!}\,\tilde V_{2k+1}(\Sig_k,\Sig_k,\Psi)
	\bigg],
	\;}
	\label{ped:eq:Sdrallorders}
\end{equation}
with $\tilde M_k:=V_{2k,\rm eff}\big|_{\Sym^k\times\Sym^k}$ and
$\Mean{\tilde u_k(\Sig_{k-1}),\Sig_k}:=\tilde V_{2k,\rm eff}(\Sig_k,\Sig_{k-1},\Psi)$
(partitions $2k=k+k$ and $2k=k+(k{-}1)+1$; every term has at most three fields). The
$k=2$ triple is exactly \eqref{p:eq:hs4:action-tilde} plus the quintic coupling ($\tilde M_2=V_4$,
$\tilde V_5=V_5$: no dressing yet); the $k=3$ even pair is \eqref{ped:eq:Sdr}
(proven); the $k=3$ odd coupling $\tilde V_7(\Sig_3,\Sig_3,\Psi)$ at
$O(\gs^5)$ is now \emph{proven too} (\cref{p:sec:hs:septic}), with the
explicit forced dressing \eqref{p:eq:hs7:dressed-vertex} --- the running rule \eqref{ped:eq:runningV}
verbatim at $n=7$; and the whole tower is now \emph{proven at every order}
(\cref{ped:sec:allorders:theorem}: the action is jointly quadratic in the
tower, the effective action is one Gaussian with an operator-continued-fraction
kernel, and closure is a solvable forced-diagonal recursion). What is uniform and established across $k$: the $1{:}2$ ratio, the
Veronese saddle $\Sig_k^*=P_k$ exactly on-shell (the collapse identities
\eqref{ped:eq:tildecollapse} pattern), homogeneous saddle-tracking gauge laws with no new
parameters, and $\IBL$ products dressed by lower
composites satisfying no new relations (\cref{ch:ibl}). The vertices grow composite as the dressing accumulates,
but the \emph{rule} \eqref{ped:eq:runningV} is one line.

\subsubsection{Scheme C (elementary, indefinite): geometric vertices plus wrong-sign
partners}
\label{ped:sec:allorders:chi}

The $\chi$-mechanism of \cref{ped:sec:milestones:chi} suggests the third bookkeeping: keep
\emph{every} vertex elementary/geometric and delegate \emph{all} subtractions to
wrong-sign partner fields. Established at $O(\gs^4)$: \eqref{ped:eq:S4chi}.
The underlying principle is general --- every composite that Route B
ever needs is the induced action of a linearly-sourced Gaussian (the Liouville
mechanism), hence can be re-localised by a partner field with the opposite-sign kernel;
what is open is the bookkeeping of the partners at $O(\gs^5)$ and beyond (the $O(\gs^5)$
feedback is now explicit, \eqref{p:eq:hs7:induced}: two strand-local chains $S_1,S_2$ in the
dressed presentation, plus the strand-split $S_3$ in the minimal one --- so either a
single $\chi$ with new couplings or a wrong-sign tower $\chi_k$ mirroring $\Sig_k$).
If it closes, this scheme is the
most geometric cubic form: all kernels are the geometric $V_n$ themselves, at
the price of an indefinite (split-signature) auxiliary sector --- which
\eqref{ped:eq:lightcone} suggests is not a bug but the signature of an enlarged gauge-fixed
theory.

